\secnumberlesssection{GLOSARIO}

{\setlength{\parskip}{0cm}

API: Interfaz de Programación de Aplicaciones (\textit{Application Programming Interface}).

ARIA: Aplicaciones Enriquecidas de Internet Accesibles (\textit{Accessible Rich Internet Applications}).

ASR: Reconocimiento Automático del Habla (\textit{Automatic Speech Recognition}).

AWS: Amazon Web Services.

CPU: Unidad Central de Procesamiento (\textit{Central Processing Unit}).

CSS: Hojas de Estilo en Cascada (\textit{Cascading Style Sheets}).

DI: Departamento de Informática.

DMAE: Degeneración Macular Asociada a la Edad.

DOM: Modelo de Objetos del Documento (\textit{Document Object Model}).

FUNDALURP: Fundación de Lucha contra la Retinitis Pigmentosa.

GB: Gigabyte.

GPU: Unidad de Procesamiento Gráfico (\textit{Graphics Processing Unit}).

HTML: Lenguaje de Marcado de Hipertexto (\textit{HyperText Markup Language}).

HTTP: Protocolo de Transferencia de Hipertexto (\textit{HyperText Transfer Protocol}).

IA: Inteligencia Artificial.

INLAB: Unidad de Inclusión Laboral (FUNDALURP).

JAWS: \textit{Job Access With Speech}.

JSON: Notación de Objetos de JavaScript (\textit{JavaScript Object Notation}).

LLM: Modelo de Lenguaje de Gran Escala (\textit{Large Language Model}).

MAE: Error Absoluto Medio (\textit{Mean Absolute Error}).

NVDA: Acceso No Visual al Escritorio (\textit{NonVisual Desktop Access}).

PS: Puntaje de Satisfacción.

RAM: Memoria de Acceso Aleatorio (\textit{Random Access Memory}).

SENADIS: Servicio Nacional de la Discapacidad.

SII: Servicio de Impuestos Internos.

SSH: Shell Segura (\textit{Secure Shell}).

SSML: Lenguaje de Marcado de Síntesis de Voz (\textit{Speech Synthesis Markup Language}).

STT: Voz a Texto (\textit{Speech to Text}), equivalente a ASR en este trabajo.

TTS: Texto a Voz (\textit{Text to Speech}), o síntesis de voz.

UC: Pontificia Universidad Católica de Chile.

UCH: Universidad de Chile.

USM: Universidad Técnica Federico Santa María (ver también UTFSM).

UTFSM: Universidad Técnica Federico Santa María.

VLM: Modelo de Lenguaje Visual (\textit{Vision-Language Model}).

VRAM: Memoria de Video (\textit{Video RAM}).

W3C: Consorcio de la World Wide Web (\textit{World Wide Web Consortium}).

WAI-ARIA: estándar del W3C para declarar la semántica de accesibilidad de elementos web (\textit{Web Accessibility Initiative -- Accessible Rich Internet Applications}).

WCAG: Pautas de Accesibilidad para el Contenido Web (\textit{Web Content Accessibility Guidelines}).
}
